% !TEX program = xelatex
% Шаблон научной статьи НИТУ «МИСиС»
% Заполните блоки ниже и скомпилируйте: xelatex article.tex 

\documentclass{misisarticle}

% --- Коды классификаторов ---
\grnti{28.23.27}
\asjc{1707, 1702}
\oecd{1.02.05}

% --- Название статьи ---
\title{Неопределённость и объяснимость в методах машинного обучения для задач компьютерного зрения и 3D-восприятия}

% --- Автор (вы) ---
\misisauthor{Соломасов Д.В.}{кафедра инженерной кибернетики, НИТУ «МИСиС»}{Москва, Россия}{m2518236@edu.misis.ru}

% --- Аннотация и ключевые слова ---
\setabstract{Работа посвящена взаимосвязи оценки неопределенности (UQ) и объяснимости (XAI) в моделях глубокого обучения для компьютерного зрения и 3D-восприятия. Различаются алеаторическая и эпистемическая неопределенность; рассматриваются методы MC-Dropout, Deep Ensembles, гетероскедастические модели, калибровка, а также LIME, SHAP, Grad-CAM, Integrated Gradients, TCAV. Особое внимание уделяется адаптации этих методов к 3D-данным: облакам точек, NeRF и монокулярной оценке глубины. Обсуждается подход Uncertainty-Aware XAI и метрики оценки качества. Цель — понять, как совместное применение UQ и XAI повышает надежность и интерпретируемость моделей 3D-восприятия.}

\setkeywords{Оценка неопределённости; Объяснимый ИИ; Компьютерное зрение; 3D-восприятие; Облака точек; Uncertainty-Aware XAI}

\begin{document}
\maketitle

\section{Введение}
Современные модели машинного обучения всё чаще применяются в критически важных областях: автономное вождение, медицинская диагностика, робототехника. Эти модели показывают высокую точность, но при этом остаются «чёрными ящиками» — мы не всегда понимаем, как они принимают решения, и не всегда можем оценить, насколько им можно доверять.

В ответ на эту проблему появились два направления исследований. Первое — оценка неопределённости — учит модель сообщать, насколько она уверена в своём ответе. Если модель не уверена, человек может перепроверить решение или запросить дополнительные данные. Второе — объяснимость — помогает понять, почему модель приняла именно такое решение, показывая, какие признаки или области изображения были для неё важны.

Особенно важны эти вопросы в задачах компьютерного зрения и 3D-восприятия. Ошибка в распознавании объекта на дороге или в медицинском снимке может стоить жизни. При этом 3D-данные (облака точек, сцены, восстановленные по фотографиям) устроены сложнее, чем обычные изображения: они не имеют чёткой структуры, могут быть разреженными и неполными. Это создаёт дополнительные трудности для оценки неопределённости и построения объяснений.

\section{Анализ работ}
В машинном обучении принято различать два типа неопределённости \cite{ref1}. Случайная (алеаторическая) неопределённость — это неустранимый шум в данных. Например, камера может снимать с шумом, объекты могут быть частично закрыты, а измерения — неточными. Даже если собрать бесконечно много данных, этот шум никуда не денется. Модельная (эпистемическая) неопределённость — это незнание модели. Она возникает, когда модель не видела похожих примеров во время обучения. 

Для оценки неопределённости на практике используются методы MC-Dropout и Deep Ensembles \cite{ref2}. Для объяснимости решений применяются методы LIME, SHAP, Grad-CAM, Integrated Gradients и другие. Однако совместное применение этих подходов (Uncertainty-Aware XAI) исследовано недостаточно, особенно для 3D-данных. В работах \cite{ref3} показано, что учёт неопределённости позволяет повысить надёжность и интерпретируемость моделей в критически важных системах.

\section{Постановка задачи}
Введём математическую постановку. Пусть $D = \{(x_i, y_i)\}_{i=1}^{N}$ — датасет, где $x_i$ — входной объект (изображение или облако точек), а $y_i$ — целевой ответ. Для задачи сегментации облаков точек $y_i \in \{1, \dots, C\}^N$ — вектор меток классов для $N$ точек.

Модель $f_\theta: X \to Y$ — это обученная нейросеть с параметрами $\theta \in \mathbb{R}^p$. Значение $f_\theta(x)$ — предсказание модели для входа $x$. Для сегментации $f_\theta(P) \in [0, 1]^{N \times C}$ — матрица вероятностей классов.

\textbf{Неопределённость.} Метод оценки неопределённости $U(f_\theta, x) \to \mathbb{R}^+$ ставит в соответствие модели и входу неотрицательное число $u = U(f_\theta, x)$, характеризующее степень неуверенности модели.

\textbf{Объяснимость.} Метод объяснимости $E(f_\theta, x) \to \mathbb{R}^{|x|}$ возвращает вектор значимости $e = E(f_\theta, x)$, где $e_j$ — значимость $j$-го элемента входа. Размерность вектора $|x|$ совпадает с числом элементов входа (пикселей или точек).

\textbf{Цель работы.} Найти функцию согласования $g: \mathbb{R}^+ \times \mathbb{R}^{|x|} \to [0, 1]$, которая связывает неопределённость и качество объяснения, а также определить правило отбора надёжных объяснений: объяснение считается надёжным, если $U(f_\theta, x) < \tau$, где $\tau$ — заданный порог.

\section{Описание методов}
Для оценки неопределённости используются следующие подходы.

\textbf{MC-Dropout.} В обычном dropout во время обучения случайно выключаются отдельные нейроны. В MC-Dropout мы оставляем dropout включённым и во время предсказаний. Прогоняем входные данные через сеть несколько раз (например, $T = 30$), каждый раз с разными случайно выключенными нейронами:
$$U_{\text{MC}}(x) = \frac{1}{T} \sum_{t=1}^{T} \left( f_{\theta_t}(x) - \bar{f}(x) \right)^2,$$
где $\theta_t$ — веса с dropout-маской на $t$-м проходе, а $\bar{f}(x)$ — среднее предсказание. Если ответы сильно различаются, модель не уверена.

\textbf{Deep Ensembles.} Обучаем не одну модель, а несколько (например, $K = 5$), каждый раз с разной случайной инициализацией весов. На этапе предсказания усредняем ответы всех моделей:
$$U_{\text{DE}}(x) = \frac{1}{K} \sum_{k=1}^{K} \left( f_{\theta_k}(x) - \bar{f}(x) \right)^2.$$
Если модели дают разные ответы — это сигнал высокой неопределённости.

Для объяснимости используются градиентные методы.

\textbf{Grad-CAM.} Показывает, какие области входа модель считает важными:
$$E_{\text{Grad-CAM}}(x) = \text{ReLU}\left( \sum_k \alpha_k A^k \right),$$
где $A^k$ — $k$-я карта признаков, $\alpha_k$ — её вес.

\textbf{Integrated Gradients.} Усредняет градиенты вдоль пути от «пустого» входа $x'$ до исходного $x$, что даёт более стабильные оценки:
$$E_{\text{IG}}(x) = (x - x') \int_0^1 \frac{\partial f_\theta(x' + \alpha(x - x'))}{\partial x} d\alpha.$$

\section*{Выводы}
В работе систематизированы методы оценки неопределённости и объяснимости для задач 3D-восприятия. Показано, что совместное применение UQ и XAI (Uncertainty-Aware XAI) позволяет не только объяснять решения моделей, но и оценивать надёжность этих объяснений. Определены направления для дальнейших экспериментов на датасетах SemanticKITTI, S3DIS, KITTI и NYU Depth v2.

% --- Список литературы (строго как в таблице) ---
\begin{thebibliography}{9}

\bibitem{ref1}
Löfström H., Löfström T., Hjort A., Yapicioglu F.R. Concerning uncertainty—a systematic survey of uncertainty-aware XAI // Philosophical Transactions of the Royal Society A. 2026. Vol. 384, No. 2327. 20250079.

\bibitem{ref2}
Gawlikowski J. et al. A survey of uncertainty in deep neural networks // Artificial Intelligence Review. 2023. Vol. 56, Suppl 1. Pp. 1513–1589.

\bibitem{ref3}
Explainability and uncertainty: Two sides of the same coin for enhancing the interpretability of deep learning models in healthcare // International Journal of Medical Informatics. 2025. Vol. 197. 105846.

\end{thebibliography}

\end{document}

\end{thebibliography}

\end{document}
